\begin{exercise}{71}
Show that\footnote{The printed identities specify no range for $n$.
Read them with $n\ge1$, $m\ge0$, and $0\le j\le n$.
Definition~1.1.37 gives $\Lambda^0_0=\varnothing$,
but $\varnothing\star\Delta^m\cong\Delta^m$.
For $n=0$, both displayed left-hand sides are
$\partial\Delta^{m+1}$; already $n=m=j=0$ gives two
vertices rather than the single vertex on the right.
The boundary cases (1) and (3) remain valid for $n=0$.}
\[
 \Lambda^n_j\star\Delta^m\cup\Delta^n\star\partial\Delta^m
 =\Lambda^{n+1+m}_j
\]
and
\[
 \partial\Delta^m\star\Delta^n\cup\Delta^m\star\Lambda^n_j
 =\Lambda^{n+1+m}_{m+1+j}.
\]
In order to do so, determine explicitly the following sub-simplicial sets of $\Delta^{n+1+m}$:
\begin{enumerate}[label=(\arabic*)]
 \item $\partial\Delta^n\star\Delta^m$
 \item $\Lambda^n_j\star\Delta^m$
 \item $\Delta^m\star\partial\Delta^n$
 \item $\Delta^m\star\Lambda^n_j$
\end{enumerate}

\begin{proof}
Put $N=n+1+m$, and let $F_t\subseteq\Delta^N$ be the
face opposite vertex $t$.
Use the concatenation identifications of
Exercise~\ref{ex:68}.
For $n\ge1$, the four answers are
\[
\begin{aligned}
 \partial\Delta^n\star\Delta^m
   &=\bigcup_{t=0}^{n}F_t,\\
 \Lambda^n_j\star\Delta^m
   &=\bigcup_{\substack{0\le t\le n\\t\ne j}}F_t,\\
 \Delta^m\star\partial\Delta^n
   &=\bigcup_{t=0}^{n}F_{m+1+t},\\
 \Delta^m\star\Lambda^n_j
   &=\bigcup_{\substack{0\le t\le n\\t\ne j}}F_{m+1+t}.
\end{aligned}
\]
Indeed, a simplex belongs to a boundary join precisely
when its image omits a vertex of the boundary factor's
block.
For a horn join, it must omit a vertex of that block
other than the horn index $j$
(Definition~1.1.37).
These descriptions include half-empty cuts: an empty
block omits every one of its vertices, and in the horn
case there is a vertex other than $j$ because $n\ge1$.
In the last two formulas the second block starts at
$m+1$, explaining the shift.

For the first requested identity, the two terms are
the unions of $F_t$ over $0\le t\le n$, $t\ne j$,
and over $n+1\le t\le N$, respectively.
Together these are exactly all faces except $F_j$,
hence $\Lambda^N_j$.
For the second identity the index sets are
$0\le t\le m$ and $t=m+1+r$ with $0\le r\le n$,
$r\ne j$.
They give all faces except $F_{m+1+j}$, as required.

When $n=0$, (1) and (2) both equal $F_0$, while
(3) and (4) both equal $F_{m+1}$, since
$\partial\Delta^0=\Lambda^0_0=\varnothing$ and join
with the empty simplicial set retains the other factor.
This also verifies the exceptional case in the footnote.
\end{proof}
\end{exercise}
